
This study examined whether keyword tagging --- the FAIR F1 findability operand --- predicts ML dataset adoption on OpenML, using negative binomial regression on a corpus of 5,217 datasets with at least one registered task.

The results support a threshold-plus-amplification characterization of the FAIR F1 mechanism:

\begin{enumerate}
\item \textbf{Binary tag presence is associated with a 22.6\% adoption advantage} (IRR=1.2263, 95\% CI [1.1681, 1.2873], p=1.$87\times 10^{-16}$, N=5,217), surviving decade fixed effects with only 12.2\% attenuation despite Cramér's V=0.823 era-tag collinearity.
\end{enumerate}

\begin{enumerate}
\item \textbf{Tag count magnitude is associated with a log-linear amplification} within tagged datasets (IRR=1.5332 per log(tag\_count+1), 95\% CI [1.4680, 1.6014], p=1.$28\times 10^{-82}$, N=2,625), with essentially zero decade confounding (attenuation ratio=1.0007).
\end{enumerate}

\begin{enumerate}
\item \textbf{Comprehensive tagging (6+ tags) is the dominant categorical tier} (IRR=1.2861, distinctly above the 1--5 tag range at p=5.$54\times 10^{-10})$, while OpenML tagging behavior is bimodal --- users tag comprehensively or not at all.
\end{enumerate}

The informative negative result (H-M3) honestly documents the limits of categorical characterization at low tag counts while refining the practical recommendation toward the 6+ tag target.

\subsubsection{7.1 Future Directions}

\textbf{Reverse causality testing:} Whether popular datasets attract tags retroactively requires tag assignment timestamps unavailable in the current corpus. A natural experiment around OpenML API version changes affecting tagging permissions would provide stronger evidence on temporal ordering.

\textbf{Direct search pathway verification:} OpenML search API log analysis (query $\rightarrow$ click $\rightarrow$ task creation sequences) would directly verify the tag-indexed search mechanism, converting proxy evidence to direct evidence.

\textbf{Multi-platform replication:} Testing FAIR F1 threshold-plus-amplification on HuggingFace, Kaggle, and UCI would examine whether the structure generalizes across repository architectures with different search mechanics.

\textbf{Hurdle model on the full corpus:} Including N\_tasks=0 datasets would estimate P(any adoption) alongside conditional adoption intensity, providing a complete picture of tagging effects across the adoption distribution.

\subsubsection{7.2 Closing}

As ML repositories grow and dataset proliferation accelerates, the structural advantages of keyword-indexed FAIR F1 compliance compound. Adding 6 or more descriptive keyword tags at dataset upload is associated with substantially higher conditional adoption on OpenML. While causal ordering requires timestamp verification beyond the available corpus, the magnitude of the association and its structural robustness across model variants support treating FAIR F1 tagging as a discoverability mechanism with measurable adoption consequences rather than a documentation checkbox.

